% s-115-p115-m005; original PDF 115, 116
\Aphorism{s-115-p115-m005}{134}
Sich nicht mehr ernst nehmen: Sache des guten Geschmacks. Gegen die feierliche Biederkeit unseres Egoismus thut uns ein tiefsinniger Carneval des Geistes mehr noth als alle Tugend-Grandezza. Was liegt an uns? Ob wir uns haben, entbehren, begeistern, opfern? Das Leben überlebt uns und nicht nur, indem es unser Gewicht überwiegt: lieber noch siegt es mit Lächeln und Feinheiten. Auf Menschengräbern wachsen bunte Blumen, mit wunderlich verzogenen Gesichtern, die spöttisch und weise lachen: kein Zweifel, dass sie uns auslachen! die Todten wie die Trauernden! und zwischen pathetischen Marmortrümmern flattert in glückseligem Augenblicks-Ernst und -Spiel ein liebetolles Schmetterlings-Pärchen\Dots{3}
\par

Alles aufrichtig, absichtslos, in flüchtiger Schönheit aus dem Weltenfüllhorn ausgeschüttet: nur der Mensch deutet mit Fingerzeigen und Zaunspfählen und lautem \Foreign{Ecce homo} auf seine Auch-Anwesenheit, nur er, der ewige Bajazzo, tritt vor die Rampe und trägt sich vor!
\par

